% =========================================================
\section{Sistema de automatización desarrollado}
\label{sec:sistema}
% =========================================================

El software se desarrolló como una aplicación de escritorio en
Python~3.11 con interfaz gráfica basada en \texttt{PySide6}~6.11, el
enlace oficial de Qt~6 para Python. La arquitectura sigue un patrón
modular en el que cada subsistema de hardware corresponde a un paquete
independiente, y la interfaz gráfica actúa como coordinador de alto
nivel.

\subsection{Arquitectura modular}
\label{sec:arquitectura}

El proyecto se organiza en torno a un paquete principal \texttt{app/}
que agrupa los siguientes módulos:

\begin{itemize}
  \item \texttt{app/laser/control\_laser.py} — gestión de la conexión
    con el láser mediante la biblioteca de enlace dinámico
    \texttt{REMOTECONTROL64.dll}. Expone una interfaz pública de lectura
    (\texttt{read\_output\_level}, \texttt{read\_eo\_delay},
    \texttt{read\_burst\_mode}, \texttt{leer\_parametros}) y de
    escritura (\texttt{set\_output\_level}, \texttt{set\_eo\_delay},
    \texttt{set\_burst\_mode}); el acceso directo a registros queda
    encapsulado y no se invoca desde la interfaz gráfica.

  \item \texttt{app/osciloscopio/control\_osciloscopio.py} —
    comunicación VXI-11 con el osciloscopio; implementa la consulta de
    parámetros de adquisición, la transferencia binaria de forma de onda
    y la lectura del preámbulo necesario para la conversión a unidades
    físicas.

  \item \texttt{app/temperatura/temperatura.py} — hilo de lectura
    continua de los cuatro sensores DS18B20 conectados al módulo
    ESP32~WROOM-32; mantiene la última lectura en memoria, indica si el
    valor es reciente y emite señales Qt cuando se actualiza.

  \item \texttt{app/medicion/medicion.py} — fachada de hardware y
    orquestador de la secuencia de medición automática; coordina el
    láser, el osciloscopio y el módulo de temperatura según el modo de
    operación seleccionado, y emite advertencias sin interrumpir la
    secuencia.

  \item \texttt{app/medicion/trigger.py} — lógica de disparo; evalúa
    la condición de inicio y fin de adquisición en función de la
    temperatura o del tiempo según el modo activo.

  \item \texttt{app/almacenamiento/almacenamiento.py} — persistencia
    de datos; guarda las muestras crudas en archivos \texttt{.npy} y los
    metadatos de cada medición en un registro \texttt{.csv} de sesión.

  \item \texttt{app/conexion/monitoreo.py} — hilo de monitoreo de
    conexión; verifica periódicamente la disponibilidad de cada
    dispositivo y actualiza el indicador de estado en la interfaz.

  \item \texttt{app/modo\_seguro/modo\_seguro.py} — procedimiento de
    seguridad descrito en la \cref{sec:auto_laser}; desactiva el láser
    de forma controlada ante una parada de emergencia, un fallo de
    comunicación o la conclusión de la secuencia.

  \item \texttt{app/exportacion/exportar\_mat.py} — conversión de una
    sesión completa al formato \texttt{.mat} para su análisis
    posterior en MATLAB.

  \item \texttt{app/gui/} — conjunto de ventanas y componentes
    gráficos.
\end{itemize}

El punto de entrada de la aplicación es \texttt{main.py}, ubicado en la
raíz del proyecto, que instancia la ventana de bienvenida y entrega el
control al ciclo de eventos de Qt.

\subsection{Decisiones de diseño}
\label{sec:decisiones}

La elección de \texttt{PySide6} frente a otras bibliotecas gráficas
responde a tres criterios: compatibilidad nativa con Qt~6, soporte
activo de largo plazo por parte de The Qt Company, y disponibilidad del
módulo \texttt{pyqtgraph} para visualización en tiempo real sin pérdida
de rendimiento.

Todas las operaciones de hardware se ejecutan en hilos secundarios
(\texttt{QThread}) independientes del hilo principal de la interfaz.
Esta separación es obligatoria en aplicaciones Qt: bloquear el hilo
principal durante una transferencia de forma de onda o una espera de
conexión provoca que la ventana quede sin respuesta. Cada hilo comunica
su resultado al hilo principal mediante señales Qt (\texttt{Signal}),
que Qt entrega de forma segura entre hilos.

El control del láser mediante \texttt{ctypes} directo, en lugar de un
subproceso separado o una capa intermedia, obedece a la restricción de
arquitectura de la \texttt{DLL}: la biblioteca
\texttt{REMOTECONTROL64.dll} es exclusivamente de 64~bits y debe
cargarse en el mismo espacio de proceso que el intérprete de Python. El
uso de \texttt{ctypes.WinDLL} desde Python~3.11 de 64~bits satisface
esta condición sin dependencias adicionales.

El osciloscopio se opera en modo de solo lectura desde la aplicación.
Las escalas vertical y horizontal, el acoplamiento, el nivel de disparo
y la longitud de registro se configuran físicamente en el instrumento y
el programa únicamente los consulta y los registra junto con cada
medición. Esta decisión responde a un criterio de trazabilidad: si el
programa pudiera modificar la configuración del instrumento, los
metadatos guardados describirían lo que el programa creyó haber
establecido y no necesariamente lo que el instrumento aplicó. Al
consultar en lugar de escribir, el osciloscopio es la única fuente de
verdad sobre su propio estado. En la \cref{sec:trabajo_futuro} se
discute la reincorporación del control de escalas como mejora
propuesta.

La automatización utiliza únicamente los canales CH1 y CH2, por lo que
las diferencias entre los modelos de la serie TDS5000B disponibles en
el laboratorio no afectan su operación. La respuesta a la consulta
\texttt{*IDN?} se guarda en el registro de cada sesión, de modo que el
instrumento empleado queda documentado en los propios datos.

\subsection{Interfaz gráfica}
\label{sec:gui}

La interfaz ofrece tres modos de operación seleccionables desde una
ventana de bienvenida: control exclusivo del láser, control exclusivo
del osciloscopio, y modo combinado que habilita la secuencia de
medición automática. Cada modo abre una ventana independiente; al
cerrarla, la aplicación regresa a la bienvenida sin necesidad de
reiniciar el proceso.

Las ventanas de control individual, mostradas en la
\cref{fig:ventanas_individuales}, sirven para la puesta a punto previa
a una secuencia: alineación del arreglo, ajuste del nivel de salida del
láser y verificación de que el osciloscopio dispara y transfiere. Ambas
comparten con la ventana principal la barra de estado de conexión y el
botón de parada de emergencia.

\begin{figure}[H]
  \centering
  \begin{subfigure}[t]{0.48\textwidth}
    \centering
    \includegraphics[width=\textwidth]{ventana_laser}
    \caption{Control del láser.}
    \label{fig:ventana_laser}
  \end{subfigure}
  \hfill
  \begin{subfigure}[t]{0.48\textwidth}
    \centering
    \includegraphics[width=\textwidth]{ventana_osciloscopio}
    \caption{Control del osciloscopio.}
    \label{fig:ventana_osciloscopio}
  \end{subfigure}
  \caption{Ventanas de control individual. La de la derecha muestra el
           bloque de solo lectura con la configuración consultada al
           instrumento, que la aplicación nunca escribe.}
  \label{fig:ventanas_individuales}
\end{figure}

La ventana principal del modo combinado
(\texttt{app/gui/ventana\_ambos.py}) se organiza en tres pestañas, que
se muestran en la \cref{fig:ventana_ambos}:

\begin{itemize}
  \item \textbf{Parámetros} — configura los parámetros del láser
    (nivel de energía, modo de disparo, longitud de ráfaga,
    \textit{EO delay}) y muestra, en modo de solo lectura, la
    configuración vigente del osciloscopio consultada al instrumento.

  \item \textbf{Medición} — expone los controles de la secuencia
    automática (modo de operación, temperatura objetivo o duración,
    botones de inicio y parada de emergencia) y muestra el estado en
    tiempo real de cada dispositivo mediante indicadores de color.

  \item \textbf{Visualizar datos} — presenta la forma de onda
    adquirida en un gráfico interactivo de \texttt{pyqtgraph} con ejes
    calibrados en tiempo (\si{\micro\second}) y voltaje
    (\si{\milli\volt}), y permite exportar la sesión a formato
    \texttt{.mat}.
\end{itemize}

\begin{figure}[H]
  \centering
  \begin{subfigure}[t]{0.9\textwidth}
    \centering
    \includegraphics[width=\textwidth]{ventana_auto_1}
    \caption{Pestaña \textit{Parámetros}.}
    \label{fig:ventana_ambos_parametros}
  \end{subfigure}

  \vspace{0.6em}

  \begin{subfigure}[t]{0.9\textwidth}
    \centering
    \includegraphics[width=\textwidth]{ventana_auto_2}
    \caption{Pestaña \textit{Medición}.}
    \label{fig:ventana_ambos_medicion}
  \end{subfigure}

  \vspace{0.6em}

  \begin{subfigure}[t]{0.9\textwidth}
    \centering
    \includegraphics[width=\textwidth]{ventana_auto_3}
    \caption{Pestaña \textit{Visualizar datos}.}
    \label{fig:ventana_ambos_visualizacion}
  \end{subfigure}
  \caption{Ventana principal del modo combinado. La barra superior
           conserva en las tres pestañas los indicadores de conexión de
           los tres dispositivos, las lecturas de los cuatro sensores de
           temperatura y el botón de parada de emergencia.}
  \label{fig:ventana_ambos}
\end{figure}

La barra superior muestra el estado de conexión de cada dispositivo
(láser, osciloscopio, ESP32) mediante indicadores de color: verde para
conexión activa, rojo para desconectado, y un botón de reconexión
disponible en todo momento. Las lecturas de los cuatro sensores de
temperatura se actualizan cada segundo en la barra superior derecha.

\subsection{Secuencia de medición automática}
\label{sec:secuencia}

El sistema implementa dos modos de medición automática. En ambos, el
láser opera de forma continua durante toda la secuencia; su activación
y desactivación forman parte del flujo de inicio y del procedimiento de
modo seguro, respectivamente.

En \textbf{modo tiempo}, el osciloscopio adquiere durante un intervalo
fijo especificado por el operador. En \textbf{modo temperatura}, el
osciloscopio actúa como integrador: inicia la adquisición
(\texttt{ACQ:STATE RUN}) cuando la temperatura de la muestra supera el
umbral objetivo en \SI{0.1}{\celsius} y la detiene
(\texttt{ACQ:STATE STOP}) cuando cae \SI{0.1}{\celsius} por debajo del
objetivo durante el enfriamiento. El número estimado de pulsos
integrados en la ventana se calcula como

\begin{equation}
  n_\text{pulsos} = (t_\text{fin} - t_\text{inicio}) \times
                    \SI{10}{\hertz}
  \label{eq:pulsos_estimados}
\end{equation}

y se almacena junto con los metadatos de la medición.

La secuencia en modo temperatura sigue los pasos descritos a
continuación:

\begin{enumerate}
  \item El operador configura los parámetros en la pestaña
    \textbf{Parámetros} y presiona \textit{Iniciar} en la pestaña
    \textbf{Medición}.
  \item \texttt{medicion.py} envía los parámetros al láser mediante la
    interfaz de \texttt{control\_laser.py} y consulta la configuración
    vigente del osciloscopio.
  \item El láser inicia el disparo.
  \item El hilo \texttt{TriggerWorker} supervisa la temperatura
    reportada por el módulo de temperatura; cuando
    $T \geq T_\text{obj} + \SI{0.1}{\celsius}$, envía
    \texttt{ACQ:STATE RUN} y registra $t_\text{inicio}$.
  \item Cuando $T \leq T_\text{obj} - \SI{0.1}{\celsius}$, envía
    \texttt{ACQ:STATE STOP} y registra $t_\text{fin}$.
  \item \texttt{medicion.py} transfiere la forma de onda mediante
    \texttt{CURVE?} junto con su preámbulo, y entrega las muestras
    crudas y los metadatos al módulo de almacenamiento. La conversión a
    unidades físicas se realiza en el análisis posterior con las
    \cref{eq:conv_voltaje,eq:conv_tiempo}.
  \item El procedimiento de modo seguro desactiva el láser.
  \item La interfaz muestra la señal adquirida y habilita el botón de
    nueva medición.
\end{enumerate}

\subsection{Formato de datos de sesión}
\label{sec:formato_datos}

Cada sesión genera una carpeta identificada por marca de tiempo que
contiene las muestras crudas de cada captura en archivos \texttt{.npy},
un registro \texttt{.csv} con los metadatos de todas las mediciones, un
archivo \texttt{metadatos\_sesion.txt} con las condiciones generales de
la sesión, y un archivo \texttt{advertencias.log} con los avisos
emitidos durante la ejecución.

La separación entre muestras crudas y metadatos es deliberada: el
archivo \texttt{.npy} conserva los valores enteros tal como los entrega
el osciloscopio, y las columnas del preámbulo (\texttt{XINCR},
\texttt{YMULT}, \texttt{YOFF}, \texttt{YZERO}, \texttt{PT\_OFF},
\texttt{NR\_PT}) permiten reconstruir la señal en unidades físicas en
cualquier momento posterior. El registro incluye además los parámetros
del láser (\texttt{output\_level}, \texttt{eo\_delay\_us},
\texttt{burst\_mode}), las escalas consultadas al osciloscopio
(\texttt{CH\_SCALE}, \texttt{HOR\_SCALE}) y los campos
\texttt{error\_flag} y \texttt{error\_desc}, que señalan si la medición
se registró en condiciones anómalas y cuál fue la anomalía.
